% Standalone mathematical section. Requires amsmath,amssymb,amsthm.
% Uses theorem, lemma, corollary, proposition, remark environments.
% Citation keys, if added by parent: BGN2018, Reiffenstein2022, DLMF.
\section{A universal law for the displacement of infinitely many zeros}
\label{sec:universal-transport}

This section supplies the convergence and regularization missing from the
earlier graph-coupled-zero construction. The basic spectral-shift mechanism
is classical; see \cite{bgn,reiffenstein}. The contribution developed here is its arithmetic asymptotic:
the total zero displacement has an explicit universal leading coefficient,
and every coefficient of its Poisson-smoothed expansion is a finite matrix
trace. No simplicity, zero-spacing, linear-independence, or pair-correlation
assumption is used.

\subsection{Inputs and a canonical measure of displacement}
\label{sec:setup}

Let $F_1,\ldots,F_n$ be phase-real primitive completed Dirichlet functions,
allowing repetitions and the pole-removed zeta completion. Thus, for a
primitive nonprincipal character $\chi$ of conductor $q$ and parity $a$,
\[
 \Lambda_\chi(s)=(q/\pi)^{(s+a)/2}\Gamma((s+a)/2)L(s,\chi),
 \qquad F_\chi(z)=\eta_\chi^{-1}\Lambda_\chi(1/2+iz),
 \qquad \eta_\chi^2=\varepsilon_\chi.
\]
For zeta set $q=1$ and
\[
 F_\zeta(z)=\xi(1/2+iz),\qquad
 \xi(s)=\tfrac12s(s-1)\pi^{-s/2}\Gamma(s/2)\zeta(s).
\]
The functional equation makes each $F_i$ real entire, including for complex
characters. Invoke the Dirichlet zero-location theorem to conclude that
all its zeros are real. Its order-one canonical product gives, off those
zeros,
\[
 g_i(z)=\frac{F_i'(z)}{F_i(z)},\qquad
 \operatorname{Im}g_i(z)<0\quad(\operatorname{Im}z>0).
\]
Indeed, with zeros listed with multiplicity,
\[
 \operatorname{Im}g_i(x+iy)
   =-y\sum_{\rho\in\mathcal Z(F_i)}\frac1{|x+iy-\rho|^2}<0.
\]
The real constant from the canonical product does not affect this formula;
the series converges locally uniformly away from the zeros.

For a fixed Hermitian $n\times n$ matrix $A$, define
\begin{align}
 P_A(z)&=\det\!\left(\operatorname{diag}(F_i(z))A+
                           \operatorname{diag}(F_i'(z))\right),
             \label{ut:defPA}\\
 M_A(z)&=A+\operatorname{diag}(g_i(z)). \label{ut:defMA}
\end{align}
The first expression is entire and real: its principal-minor expansion
has real coefficients multiplying products of real entire functions.
Away from the input zeros,
$P_A=(\prod_iF_i)\det M_A$. Since $\operatorname{Im}M_A\prec0$ in the
upper half-plane, $M_A$ is invertible there. Consequently $P_A$ is nonzero
and all its zeros are real.

Fix $0\ne B\succeq0$ Hermitian, of rank $r$, and set $A_1=A+B$.
At real $x$ avoiding all input and output zeros, define
\begin{equation}\label{ut:defshift}
 \sigma_{A,B}(x)=n_+(M_{A+B}(x))-n_+(M_A(x)).
\end{equation}
Values on the excluded discrete set are immaterial for integration.
Here $n_+$ counts strictly positive eigenvalues with multiplicity.
Since $B\succeq0$ and has rank $r$,
\begin{equation}\label{ut:shiftbound}
       0\leq\sigma_{A,B}(x)\leq r.
\end{equation}
This is the canonical rightward zero-crossing count. In particular, it
avoids the ambiguous subtraction of two infinite zero-counting functions.

\begin{theorem}[Exact displacement measure and trace identity]
\label{ut:exact}
There is a holomorphic logarithm $\mathcal L_{A,B}$ on the upper
half-plane satisfying
\[
 e^{\mathcal L_{A,B}(z)}=\frac{P_{A+B}(z)}{P_A(z)},\qquad
 0<\operatorname{Im}\mathcal L_{A,B}(z)<r\pi,
 \qquad \mathcal L_{A,B}(iy)\longrightarrow0\quad(y\to\infty).
\]
Its boundary imaginary part is $\pi\sigma_{A,B}$ almost everywhere, and
\begin{equation}\label{ut:poisson}
 \boxed{\quad
 \int_{\mathbb R}\frac{y^2\sigma_{A,B}(t)}{t^2+y^2}\,dt
       =y\operatorname{Im}\mathcal L_{A,B}(iy)
       \quad(y>0).\quad}
\end{equation}
More generally,
\begin{equation}\label{ut:poisson-general}
 \operatorname{Im}\mathcal L_{A,B}(x+iy)
     =\int_{\mathbb R}\frac{y\sigma_{A,B}(t)}{(t-x)^2+y^2}\,dt.
\end{equation}
Writing $\nu_A$ for the locally finite zero-counting measure of $P_A$,
including multiplicity, one has the distributional equality
\begin{equation}\label{ut:distribution}
         d\sigma_{A,B}=d\nu_A-d\nu_{A+B}.
\end{equation}
Consequently every $\varphi\in C_c^1(\mathbb R)$ satisfies the exact,
finite identity
\begin{equation}\label{ut:testtrace}
 \sum_{P_{A+B}(\beta)=0}\!\operatorname{ord}_\beta(P_{A+B})\varphi(\beta)
 -\sum_{P_A(\alpha)=0}\!\operatorname{ord}_\alpha(P_A)\varphi(\alpha)
   =\int_{\mathbb R}\varphi'(t)\sigma_{A,B}(t)\,dt.
\end{equation}
Common zeros cancel with their exact multiplicities.
\end{theorem}

\begin{proof}
Write $B=VV^*$ with $V$ an $n\times r$ matrix of full column rank. Set
\[
 H(z)=V^*M_A(z)^{-1}V.
\]
For every invertible matrix $M$,
\[
 \operatorname{Im}(M^{-1})
       =-M^{-*}(\operatorname{Im}M)M^{-1}.
\]
Thus $\operatorname{Im}H(z)\succ0$ in the upper half-plane. The eigenvalues
of $I_r+H(z)$ all lie in that half-plane, so its principal matrix logarithm
is holomorphic. Put
\[
 \mathcal L_{A,B}(z)=\operatorname{tr}\log(I_r+H(z)).
\]
The sum of the principal arguments of its $r$ eigenvalues lies strictly
between $0$ and $r\pi$. The determinant lemma gives
$e^{\mathcal L_{A,B}}=\det M_{A+B}/\det M_A=P_{A+B}/P_A$.
The asymptotic lemma below gives $H(iy)=O(1/\log y)$ and hence the asserted
normalization at infinity.

At a real regular point $x$, the matrices $M_A(x)$ and $H(x)$ are Hermitian.
Approach from the upper half-plane. Each negative eigenvalue of $I+H(x)$
has limiting argument $\pi$ and each positive eigenvalue has limiting
argument zero. Therefore
\[
 \pi^{-1}\operatorname{Im}\mathcal L_{A,B}(x+i0)=n_-(I_r+H(x)).
\]
To identify this with \eqref{ut:defshift}, take the inertia of
\[
 \begin{pmatrix}M_A(x)&V\\ V^*&-I_r\end{pmatrix}
\]
using its two Schur complements. They give, respectively,
$n_+(M_A(x))+n_-(I_r+H(x))$ and $n_+(M_A(x)+B)$.
This proves the boundary formula.

The imaginary part of $\mathcal L_{A,B}$ is a bounded harmonic function
on the upper half-plane. The Poisson representation for bounded harmonic
functions gives \eqref{ut:poisson-general}. There is no additional linear
term or singular boundary measure: bounded harmonic functions are the
Poisson extensions of their bounded almost-everywhere boundary values.
This proves \eqref{ut:poisson} as well; its integral is absolutely finite
by \eqref{ut:shiftbound}.

For completeness, multiplicities in \eqref{ut:distribution} can be read
directly from the entire quotient, without a simplicity assumption.
At a real $x_0$, locally write
\[
 \frac{P_{A+B}(z)}{P_A(z)}=(z-x_0)^k a(z),\qquad
 k=\operatorname{ord}_{x_0}P_{A+B}-\operatorname{ord}_{x_0}P_A,
\]
with $a$ analytic, real and nonzero near $x_0$. The boundary argument
of $\log(z-x_0)$ drops by $\pi$ when crossing $x_0$ from left to right.
Hence $\sigma_{A,B}$ jumps by $-k$. Away from this discrete divisor it
is locally constant. These are precisely the atoms of
$\nu_A-\nu_{A+B}$. Integration by parts gives \eqref{ut:testtrace}.
\end{proof}

\begin{remark}[Why this measures motion rather than just a count]
For a rank-one addition, $1+H(z)$ is a scalar meromorphic function with
positive imaginary part. Its real poles and zeros interlace after common
factors are canceled. The set where its boundary value is negative is a
disjoint union of intervals from an old zero to the following new zero;
$\sigma_{A,B}$ is their indicator. Thus its integral over a window is
the total length of the moving-zero intervals clipped to that window.
For general $B$, decompose it into $r$ positive rank-one terms and apply
this observation successively. The crossing counts add and telescope
to \eqref{ut:defshift}, so their total is independent of this decomposition.
No separately divergent sum $\sum\beta-\sum\alpha$ is used.
\end{remark}

\subsection{Arithmetic input at large complex arguments}

Put
\begin{equation}\label{ut:conductorconstants}
 c_i=\tfrac12\log\!\left(\frac{q_i}{2\pi}\right),\qquad
 C=\operatorname{diag}(c_1,\ldots,c_n),\qquad
 \ell_y=\tfrac12\log y.
\end{equation}

\begin{lemma}[Uniform radial logarithmic derivative]\label{ut:radial}
For every compact subset $K$ of the upper half-plane,
\begin{equation}\label{ut:radial-asymptotic}
 g_i(Tz)=-i\left(\tfrac12\log T+c_i+\tfrac12\log(-iz)\right)
                    +O_K(T^{-1})\qquad(z\in K).
\end{equation}
The logarithm is the principal one; $-iz$ is in the right half-plane.
In particular
\begin{equation}\label{ut:imag-asymptotic}
 g_i(iy)=-i(\ell_y+c_i)+O(y^{-1}).
\end{equation}
All conductors, characters, and matrices are fixed as $T$ or $y$ tends
to infinity.
\end{lemma}

\begin{proof}
For a primitive nonprincipal character the functional equation gives
\[
 g_i(Tz)=-i\frac{\Lambda_{\overline\chi}'(1/2-iTz)}
                         {\Lambda_{\overline\chi}(1/2-iTz)}.
\]
Write $s=1/2-iTz$. Uniformly for $z\in K$, $\operatorname{Re}s$ tends
linearly to $+\infty$ and $s$ stays in a closed subsector of the right
half-plane. Logarithmic differentiation of the completion gives
\[
 \frac{\Lambda_{\overline\chi}'(s)}{\Lambda_{\overline\chi}(s)}
 =\tfrac12\log(q_i/\pi)
       +\tfrac12\psi((s+a_i)/2)
       +\frac{L'(s,\overline\chi)}{L(s,\overline\chi)}.
\]
The functional equation and completion are standard \cite{dlmfdirichlet}.
The absolutely convergent Euler-product logarithmic derivative is
exponentially small in $\operatorname{Re}s$. The uniform Stirling
expansion $\psi(w)=\log w+O(|w|^{-1})$ gives
\[
 \tfrac12\log T+c_i+\tfrac12\log(-iz)+O_K(T^{-1}).
\]
For zeta, logarithmic differentiation of the pole-removing factor adds
$1/s+1/(s-1)=O_K(T^{-1})$, so the same formula holds with $q_i=1$.
\end{proof}

\subsection{The macroscopic transport law}

\begin{theorem}[Universal macroscopic zero transport]\label{ut:main}
With the preceding fixed inputs, matrices, and $0\ne B\succeq0$, the
positive measures
\begin{equation}\label{ut:weaklimit}
   (\log T)\sigma_{A,B}(Tu)\,du
        \ \longrightarrow\ \frac{2\operatorname{tr}B}{\pi}\,du
        \qquad(T\to\infty)
\end{equation}
converge against every continuous compactly supported test function on
$\mathbb R$. Consequently, for every fixed pair $a<b$,
\begin{equation}\label{ut:windowlaw}
 \boxed{\quad
 \int_{aT}^{bT}\sigma_{A,B}(x)\,dx
       \sim \frac{2\operatorname{tr}B}{\pi}(b-a)
                          \frac{T}{\log T}.
 \quad}
\end{equation}
In particular,
\begin{equation}\label{ut:symmetriclaw}
 \int_{-T}^{T}\sigma_{A,B}(x)\,dx
       \sim\frac{4\operatorname{tr}B}{\pi}\frac{T}{\log T}.
\end{equation}
For a graph edge of added conductance $h>0$, take
$B=h(e_a-e_b)(e_a-e_b)^*$; then $\operatorname{tr}B=2h$, and the
coefficient in \eqref{ut:symmetriclaw} is $8h/\pi$.
Thus the unregularized total positive displacement is infinite, with a
precise growth law. Its leading coefficient depends on neither the
initial graph nor the characters, conductors, or individual zero spacings.
\end{theorem}

\begin{proof}
Write $b=\operatorname{tr}B>0$. Lemma~\ref{ut:radial} gives, uniformly on
compact subsets of the upper half-plane,
\[
 M_A(Tz)=-\tfrac i2(\log T)I+O(1),\qquad
 M_A(Tz)^{-1}=\frac{2i}{\log T}I+O((\log T)^{-2}).
\]
The defining principal logarithm therefore yields
\begin{equation}\label{ut:harmonic-limit}
 (\log T)\operatorname{Im}\mathcal L_{A,B}(Tz)
                       \longrightarrow 2b.
\end{equation}
Let $d\mu_T(u)=(\log T)\sigma_{A,B}(Tu)\,du$. Changing variables in
\eqref{ut:poisson-general} shows that, for every $x\in\mathbb R$, $y>0$,
\begin{equation}\label{ut:scaled-poisson}
 \int_{\mathbb R}\frac{y}{(u-x)^2+y^2}\,d\mu_T(u)
        =(\log T)\operatorname{Im}\mathcal L_{A,B}(T(x+iy))
        \longrightarrow 2b.
\end{equation}
We justify passage from these positive transforms to measures explicitly.
For fixed $R>0$, the kernel $R/(u^2+R^2)$ is at least $1/(2R)$ on
$[-R,R]$. Hence the masses $\mu_T([-R,R])$ are uniformly bounded for
all sufficiently large $T$. Every sequence $T\to\infty$ has a vaguely
convergent subsequence, with positive locally finite limit $\mu$.

The Poisson transforms also have uniformly negligible tails in the
required iterated-limit sense. Fix $x,y$ and take $R>2|x|+2y$.
For $|u|\ge R$,
\[
 \frac{y}{(u-x)^2+y^2}
       \leq\frac{4y}{u^2}
       \leq\frac{8y}{R}\frac{R}{u^2+R^2}.
\]
Thus \eqref{ut:scaled-poisson}, evaluated at $iR$, gives
\[
 \limsup_{T\to\infty}
 \int_{|u|\ge R}\frac{y}{(u-x)^2+y^2}\,d\mu_T(u)
             \leq\frac{16by}{R}.
\]
It follows, first using compact cutoffs and then letting $R\to\infty$,
that every subsequential limit satisfies
\[
 \int_{\mathbb R}\frac{y}{(u-x)^2+y^2}\,d\mu(u)=2b
                       \qquad(x\in\mathbb R, y>0).
\]
Uniqueness of the Poisson representation gives
$d\mu=(2b/\pi)\,du$. One may see this without a representation theorem
by integrating against a smooth compactly supported function of $x$ and
letting $y\downarrow0$; the normalized Poisson kernels form an approximate
identity, giving $\int\varphi\,d\mu=(2b/\pi)\int\varphi$.
Since every subsequential limit is this same measure,
\eqref{ut:weaklimit} follows. The limiting measure has no atoms at $a,b$,
so compact-interval masses converge as well. Changing variables once more
proves \eqref{ut:windowlaw}; the other assertions follow immediately.
\end{proof}

\subsection{All orders, including the first conductor and graph corrections}

Define the finite Poisson-smoothed displacement
\[
 \mathfrak D_{A,B}(y)
       =\int_{\mathbb R}\frac{y^2\sigma_{A,B}(t)}{t^2+y^2}\,dt.
\]

\begin{theorem}[Complete expansion]\label{ut:expansion}
For every fixed positive integer $m$,
\begin{equation}\label{ut:allorders}
 \boxed{\quad
 \frac{\mathfrak D_{A,B}(y)}y
   =\sum_{k=1}^{m}\frac{(-1)^{k+1}}{k\ell_y^k}
        \operatorname{Im}\operatorname{tr}
        \left((C+i(A+B))^k-(C+iA)^k\right)
        +O(\ell_y^{-m-1}).\quad}
\end{equation}
In particular,
\begin{align}
 \frac{\mathfrak D_{A,B}(y)}y
  &=\frac{\operatorname{tr}B}{\ell_y}
       -\frac{\operatorname{tr}(CB)}{\ell_y^2}
       +\frac{\kappa_3(A,B,C)}{\ell_y^3}
       +O(\ell_y^{-4}),\label{ut:threeorders}\\
 \kappa_3(A,B,C)
  &=\operatorname{tr}(C^2B)
       -\frac13\operatorname{tr}\left((A+B)^3-A^3\right)\notag\\
  &=\operatorname{tr}(C^2B)-\operatorname{tr}(A^2B)
       -\operatorname{tr}(AB^2)-\tfrac13\operatorname{tr}(B^3).
       \label{ut:thirdcoefficient}
\end{align}
The matrices need not commute. The first correction records conductor
data; the initial coupling first appears in the third displayed order.
\end{theorem}

\begin{proof}
By \eqref{ut:imag-asymptotic},
\[
 M_A(iy)=-i\ell_y\left(I+\frac{C+iA}{\ell_y}
                                      +O((y\ell_y)^{-1})\right).
\]
Use the convergent matrix logarithm series at $I$ for sufficiently large
$y$, take traces, and subtract the expressions for $A+B$ and $A$.
The factors $(-i\ell_y)^n$ cancel. Both logarithms use the branch tending
to zero after this cancellation, agreeing with $\mathcal L_{A,B}(iy)$.
The error from the $O(y^{-1})$ arithmetic remainder is
$O((y\ell_y)^{-1})$, smaller than every fixed negative power of $\ell_y$.
The usual power-series remainder gives \eqref{ut:allorders}.

For the first three terms use cyclicity of trace:
\[
 \operatorname{Im}\operatorname{tr}(C+iA)=\operatorname{tr}A,
\]
\[
 \operatorname{Im}\operatorname{tr}(C+iA)^2=2\operatorname{tr}(CA),
 \quad
 \operatorname{Im}\operatorname{tr}(C+iA)^3
             =3\operatorname{tr}(C^2A)-\operatorname{tr}(A^3).
\]
All displayed traces of Hermitian products are real. Expanding the cubic
difference under the trace proves \eqref{ut:thirdcoefficient}.
\end{proof}

\begin{corollary}[Explicit single-edge correction]
For $A=L$ a graph Laplacian and $B=hvv^*$ with $v=e_a-e_b$,
\begin{align*}
 \frac{\mathfrak D_{L,B}(y)}y
 ={}&\frac{2h}{\ell_y}-\frac{h(c_a+c_b)}{\ell_y^2}\\
 &+\frac{h(c_a^2+c_b^2)-h\|Lv\|^2
            -2h^2v^*Lv-\frac83h^3}{\ell_y^3}
       +O(\ell_y^{-4}).
\end{align*}
\end{corollary}

\subsection{An exactly computable two-vertex example}

Place the same zeta input $F(z)=\xi(1/2+iz)$ at two vertices joined by
an edge of conductance $w>0$. Then
\[
 P_w(z)=F'(z)\bigl(F'(z)+2wF(z)\bigr).
\]
Increasing the conductance by $h>0$ leaves the $F'$ factor unchanged.
Put
\[
 a(y)=\frac{\xi'(1/2+y)}{\xi(1/2+y)}>0\qquad(y>0).
\]
Its positivity also follows from the real-zero product. Functional
symmetry gives $g(iy)=-ia(y)$, and the exact determinant identity reduces
to the explicit formula
\begin{equation}\label{ut:twovertex}
 \boxed{\quad
 \mathfrak D_{w,h}(y)
  =y\left[\arctan\!\left(\frac{2(w+h)}{a(y)}\right)
             -\arctan\!\left(\frac{2w}{a(y)}\right)\right].
 \quad}
\end{equation}
For $y>1/2$, one can evaluate this without locating any zeros, using
\[
 a(y)=\frac1s+\frac1{s-1}-\frac12\log\pi
           +\frac12\psi(s/2)+\frac{\zeta'(s)}{\zeta(s)},
          \qquad s=y+\tfrac12.
\]
Equation \eqref{ut:twovertex} is exact. The associated unsmoothed
transport obeys
\[
 \int_{-T}^{T}\sigma_{w,h}(x)\,dx
                  \sim\frac{8h}{\pi}\frac{T}{\log T}.
\]
This example provides an independent sign and normalization check for
both universal constants and all the matrix-trace coefficients.

\paragraph{Precise scope.}
The results concern the explicitly constructed entire determinants.
They do not identify a new Euler product or an automorphic $L$-function.
The zero-location theorem is used solely to supply positive spectral
shift; the asymptotic at imaginary arguments follows independently from
the functional equation, the gamma factor, and absolute convergence on
the right. All limits fix the conductors and the coupling matrices.
The proofs do not imply a microscopic spacing law, and the smoothed
higher-order expansion is not asserted term by term for the unsmoothed
window integral.
